\documentclass[11p]{article}
% Packages
\usepackage{amsmath}
\usepackage{graphicx}
\usepackage[swedish]{babel}
\usepackage[
    backend=biber,
    style=authoryear-ibid,
    sorting=ynt
]{biblatex}
\usepackage[utf8]{inputenc}
\usepackage[T1]{fontenc}
\usepackage{textcomp}
\usepackage{hyperref}
%Källor
\addbibresource{mall.bib}
\graphicspath{ {./images/} }

\title{Labboration kaströrelse \\ \small Fysik 1}
\author{Alvin Högdal}
\date{\today}

\begin{document}

    \begin{titlepage}
        \begin{center}
            \vspace*{1cm}

            \Huge
            \textbf{Labboration kaströrelse}

            \vspace{0.5cm}
            \LARGE
            Fysik 2

            \vspace{1.5cm}

            \textbf{Alvin Högdal}

            \vfill


            \vspace{0.8cm}

            \includegraphics[width=0.4\textwidth]{../images/NTI Gymnasiet_Symbol_print_svart.png}

            \Large
            Teknikprogrammet\\
            NTI Gymnasiet\\
            Umeå\\
            \today

        \end{center}
    \end{titlepage}
% Om arbetet är långt har det en innehållsförteckning, annars kan den utelämnas
    \tableofcontents
    \newpage
    \section{Syfte och frågeställning}
    Att beräkna utgångshastigheten för kulkanonen och därefter kunna beräkna vart kulan kommer landa på bordet och på golvet
    \section{Material och metod}

    \subsection{Material:}
    rada upp dina frågor i punktform
    \begin{itemize}
        \item Kulkanon
        \item Kula
        \item Måttband
        \item kulmål
    \end{itemize}

    \subsection{Metod:}
    Beräkna utgångshastigheten för kulkanonen enligt formeln för stighöjd
    \begin{equation}
        Y_{max} = \frac{v_0^2*\sin^2(\alpha)}{2g}
    \end{equation}
    genom att kolla hut högt kulan far om utgångsvinkeln \alpha \text{ är 90} $^{\circ}$. \\
    Använd sedan det resultatet för att beräkna vart kulan hamnar om utgångsvinkeln \alpha \text{ är 60} $^{\circ}$.
    Beräkna också vart kulan hamnar om den träffar golvet.
    \paragraph{}

    \section{Resultat}
    Kulan flyger 129 cm upp i luften
    Om man skriver om formeln för stighöjd så att den ger resultatet för hastigheten istället för maxhöjd
    \begin{equation}
       v_0 = \sqrt(\frac{2g*Y_{max}}{\sin^2(\alpha)})
    \end{equation}
    Detta ger resultatet {v_0} = 5,03 m/s.

    Nu kan man beräkna kastvidden enligt formeln för kastvidd
    \begin{equation}
        X_{max} = \frac{v_0^2*\sin^2(\alpha)}{g}
    \end{equation}

    när \alpha \text{ är 60} $^{\circ}$.

    Detta ger resultatet 2.23 m.

    Avståndet till golvet är 73 cm
    Man kan beräkna när kulan träffar golvet genom att ställa upp en funktion för kaströrelsen och kolla när den skär linjen som representerar golvet.
    Funktionen för kulans rörelse i Y led är
    \begin{equation}
        Y = v_0*\sin(\alpha) * T - \frac{(g*T^2)}{2}
    \end{equation}


    \section{Analys}




\end{document}
